\documentclass[11pt,a4paper]{article}

% --- Packages ---
\usepackage[margin=.75in]{geometry}
\usepackage[utf8]{inputenc}
\usepackage[T1]{fontenc}
\usepackage{lmodern} % Fix for font not found error
\usepackage{titlesec}
\usepackage{enumitem}
\usepackage{hyperref}
\usepackage{xcolor}
\usepackage{fancyhdr}
\usepackage{comment}
\usepackage{cleveref}
\usepackage[bottom]{footmisc}
\usepackage{graphicx}
\usepackage{longtable}
\usepackage{booktabs}
\usepackage{array}
\usepackage{calc}
\usepackage{etoolbox}

% Standard Pandoc setup
\providecommand{\tightlist}{%
  \setlength{\itemsep}{0pt}\setlength{\parskip}{0pt}}

\crefformat{section}{\S#2#1#3}
\crefformat{subsection}{\S#2#1#3}
\crefformat{subsubsection}{\S#2#1#3}
\crefformat{figure}{#2Figure~#1#3}
\crefformat{footnote}{#2\footnotemark[#1]#3}

\linespread{0.9} % Riduce lo spazio tra le linee

% --- Colors ---
\definecolor{primary}{RGB}{0, 0, 0}
\definecolor{links}{HTML}{0000EE}
\definecolor{blue}{HTML}{26428B}

\hypersetup{
    colorlinks=true,
    linkcolor=links,
    urlcolor=links,
    citecolor=links,
}

% --- Section Formatting ---
\titleformat{\section}
  {\normalfont\Large\bfseries\color{blue}}{\thesection}{1em}{}
\titleformat{\subsection}
{\normalfont\large\bfseries\color{blue}}{\thesubsection}{1em}{}
\titleformat{\subsubsection}
{\normalfont\bfseries\color{blue}}{\thesubsubsection}{1em}{}

% --- Header and Footer ---
\pagestyle{fancy}
\fancyhead[L]{\small Giuseppe Capasso}
\fancyhead[R]{\small intro-to-linux}
\fancyfoot[C]{\thepage}
\renewcommand{\headrulewidth}{0.4pt}

% --- Custom Title Command ---
\newcommand{\proposaltitle}[1]{
    \begin{center}
        \vspace*{0.1cm}
        {\LARGE \bfseries \color{blue} #1} \\[1.5ex]
        {\large \textbf{Giuseppe Capasso}} \\[1ex]
                \small{
            \vspace{0.1cm}
            \textbf{Materia:} intro-to-linux\\
        }
            \end{center}
    \vspace{0.3cm}
    \hrule height 0.5pt
    \vspace{0.5cm}
}

\begin{document}

\proposaltitle{Mount point Backup}

Laboratorio Pratico Lezione 3: Storage \& Backup Management Scenario:
Sei il nuovo Junior SysAdmin della ``ITS Tech Solutions''. Il CTO ti ha
assegnato un nuovo server (la tua VM) a cui è stato appena collegato un
nuovo disco fisico. Il tuo compito è configurarlo per ospitare i dati
aziendali e impostare un sistema di backup automatico. Tempo stimato: 90
minuti Prerequisiti: Macchina Virtuale spenta (inizialmente).

Parte 1: Hardware Provisioning (VirtualBox) Obiettivo: Simulare
l'inserimento fisico di un disco nel server. 1. Spegni la VM se è
accesa. 2. Vai su Impostazioni -\textgreater{} Archiviazione. 3.
Seleziona il Controller: SATA. 4. Clicca sull'icona ``Aggiungi disco
rigido'' (icona disco con il +). 5. Seleziona Crea -\textgreater{} VDI
-\textgreater{} Allocato dinamicamente. 6. Imposta la dimensione a 4 GB
(o 5GB). Chiamalo Disco\_Dati. 7. Conferma e avvia la VM.

Parte 2: Partizionamento (fdisk/cfdisk) Obiettivo: Suddividere il disco
per scopi diversi. 1. Apri il terminale e identifica il nuovo disco:
Bash lsblk

(Dovresti vedere un device, probabilmente sdb, da 4GB senza partizioni).
2. Lancia l'utility di partizionamento: Bash sudo cfdisk /dev/sdb

(Se ti chiede il tipo di etichetta, scegli gpt). 3. Crea la seguente
struttura:

• Partizione 1: 2 GB (Sarà usata per i ``Dati Produzione''). •
Partizione 2: 2 GB (Spazio rimanente, sarà usata per i ``Backup
Locali''). 4. Seleziona {[} Write {]}, digita yes per confermare e poi
{[} Quit {]}. 5. Verifica che le partizioni esistano: Bash lsblk

(Ora dovresti vedere sdb1 e sdb2).

Parte 3: Filesystem e Formattazione Obiettivo: Rendere le partizioni
utilizzabili. 1. Formatta entrambe le partizioni con il filesystem
standard ext4: Bash sudo mkfs.ext4 /dev/sdb1 sudo mkfs.ext4 /dev/sdb2

Parte 4: Mounting Persistente (/etc/fstab) Obiettivo: Configurare il
server affinché i dischi siano pronti ad ogni riavvio. 1. Crea i punti
di mount (le cartelle vuote): Bash sudo mkdir -p /mnt/produzione sudo
mkdir -p /mnt/backup\_drive

\begin{enumerate}
\def\labelenumi{\arabic{enumi}.}
\setcounter{enumi}{1}
\item
  Ottieni gli UUID delle partizioni (copiali su un blocco note o tienili
  pronti): Bash sudo blkid
\item
  Modifica il file /etc/fstab per rendere il mount permanente: Bash sudo
  nano /etc/fstab
\item
  Aggiungi in fondo al file queste due righe (sostituisci
  TUO\_UUID\_\ldots{} con i codici reali): Plaintext
\end{enumerate}

\# Mount point Dati Produzione UUID=inserisci-uuid-sdb1

/mnt/produzione

ext4

defaults

0

2

/mnt/backup\_drive

ext4

defaults

0

2

\section{Mount point Backup}\label{mount-point-backup}

UUID=inserisci-uuid-sdb2

\begin{enumerate}
\def\labelenumi{\arabic{enumi}.}
\setcounter{enumi}{4}
\tightlist
\item
  TEST CRUCIALE: Prima di riavviare, verifica che non ci siano errori!
  Bash sudo mount -a
\end{enumerate}

(Se il comando non restituisce output, è andato tutto bene. Se dà
errore, correggi fstab o la VM non partirà). 6. Verifica che siano
montati: Bash df -h

\begin{enumerate}
\def\labelenumi{\arabic{enumi}.}
\setcounter{enumi}{6}
\tightlist
\item
  Gestione Permessi: Di default, root possiede queste cartelle.
  Assegnale al tuo utente: Bash sudo chown -R \(USER:\)USER
  /mnt/produzione sudo chown -R \(USER:\)USER /mnt/backup\_drive
\end{enumerate}

Parte 5: Simulazione Dati e Backup ``Cold'' (TAR) Obiettivo: Creare un
archivio compresso di salvataggio. 1. Generazione Dati: Crea una
struttura di file finti in produzione: Bash mkdir
/mnt/produzione/progetto\_web echo ``Codice importante'' \textgreater{}
/mnt/produzione/progetto\_web/index.html echo ``Configurazione segreta''
\textgreater{} /mnt/produzione/progetto\_web/config.php

\begin{enumerate}
\def\labelenumi{\arabic{enumi}.}
\setcounter{enumi}{1}
\item
  Esecuzione Backup (Full): Crea un archivio .tar.gz di tutto il
  progetto e salvalo nel disco di backup. Bash tar -czvf
  /mnt/backup\_drive/backup\_full\_v1.tar.gz
  /mnt/produzione/progetto\_web
\item
  Simulazione Disastro: Un utente distratto cancella tutto!
\end{enumerate}

Bash rm -rf /mnt/produzione/progetto\_web

\begin{enumerate}
\def\labelenumi{\arabic{enumi}.}
\setcounter{enumi}{3}
\tightlist
\item
  Restore (Ripristino): Recupera i dati dall'archivio. Bash \#
  Attenzione alla posizione: tar scompatta dove gli dici o nel percorso
  assoluto se registrato tar -xzvf
  /mnt/backup\_drive/backup\_full\_v1.tar.gz -C /
\end{enumerate}

(Nota: Abbiamo usato -C / perché il tar è stato creato con percorsi
assoluti. Verifica se i file sono tornati in /mnt/produzione).

Parte 6: Backup Incrementale e Sync (RSYNC) Obiettivo: Mantenere una
copia speculare (mirror) sempre aggiornata. 1. Crea una cartella per il
mirror sul disco di backup: Bash mkdir
/mnt/backup\_drive/mirror\_giornaliero

\begin{enumerate}
\def\labelenumi{\arabic{enumi}.}
\setcounter{enumi}{1}
\tightlist
\item
  Lancia la prima sincronizzazione: Bash rsync -av /mnt/produzione/
  /mnt/backup\_drive/mirror\_giornaliero/
\end{enumerate}

(Nota lo slash finale / dopo produzione: significa ``il contenuto di'',
non la cartella stessa). 3. Modifica Dati: Aggiungi un nuovo file in
produzione. Bash touch /mnt/produzione/progetto\_web/nuovo\_file.txt

\begin{enumerate}
\def\labelenumi{\arabic{enumi}.}
\setcounter{enumi}{3}
\tightlist
\item
  Sync Incrementale: Rilancia lo stesso comando rsync di prima. •
  Osserva l'output: copierà solo nuovo\_file.txt. Non ricopia tutto il
  resto. Questo è il vantaggio di rsync!
\end{enumerate}

Challenge Finale (Per chi finisce prima) Obiettivo: Automatizzare il
backup con Cron. Configura il sistema affinché esegua il comando rsync
(quello della Parte 6) automaticamente ogni minuto (solo per test
didattico).

1. Apri il crontab: crontab -e 2. Aggiungi la riga magica (cerca su
Google ``Crontab every minute syntax'' se non la ricordi): \emph{ } * *
* rsync -av /mnt/produzione/ /mnt/backup\_drive/mirror\_giornaliero/
\textgreater\textgreater{} /tmp/backup.log 3. Crea un file in
produzione, aspetta 60 secondi e controlla se appare magicamente nel
backup.

Troubleshooting Rapido • Errore ``Permission Denied'': Hai dimenticato
il sudo chown nella Parte 4. Le cartelle montate appartengono a root
finché non le cambi. • Errore ``Device is busy'' (quando smonti): Sei
ancora dentro la cartella col terminale. Fai cd \textasciitilde{} e
riprova. • La VM va in Emergency Mode al riavvio: Hai sbagliato l'UUID
in /etc/fstab. • Soluzione: Inserisci la password di root, dai nano
/etc/fstab, commenta (metti \#) davanti alle righe aggiunte, salva e
riavvia. Poi riprova con calma.



\end{document}
